\chapter{Congenital anomalies (birth differences)}

\section*{Definition and Overview}
Congenital anomalies, also referred to as birth differences, birth defects, or congenital disorders, are structural, functional, or metabolic abnormalities present at birth. These anomalies can arise during intrauterine development and may be identified prenatally, at birth, or later in infancy and childhood. They range in severity from minor structural variations with no clinical significance to major malformations that cause chronic disability or are incompatible with life. Congenital anomalies are a leading cause of infant mortality and paediatric morbidity globally, requiring multidisciplinary care and individualised long-term management.

\section*{Epidemiology}
Congenital anomalies affect approximately 2\% to 3\% of all births in Australia, which is consistent with global estimates of around 6\% of births worldwide when including disorders diagnosed later in life. Major structural anomalies, such as congenital heart defects, neural tube defects, and Down syndrome, represent a significant proportion of cases. The prevalence of specific anomalies varies by maternal age (particularly for chromosomal trisomies), genetic background, geographic location, and socioeconomic status. In Australia, the implementation of nationwide screening programmes and folate fortification of flour has significantly reduced the incidence of certain anomalies, such as neural tube defects, though they remain a major public health focus.

\section*{Aetiology and Pathophysiology}
The development of congenital anomalies is multifactorial, involving genetic, environmental, and socio-demographic factors, though a significant proportion remain idiopathic.
\begin{itemize}
    \item \textbf{Genetic mutations and chromosomal abnormalities:} Aneuploidies (e.g.\ Down syndrome/Trisomy 21, Edwards syndrome/Trisomy 18), single-gene mutations (e.g.\ cystic fibrosis, achondroplasia), or microdeletions/duplications.
    \item \textbf{Teratogenic exposures:} Maternal ingestion of harmful substances during critical periods of organogenesis (weeks 3 to 8 of gestation), such as alcohol (causing Foetal Alcohol Spectrum Disorder), tobacco, recreational drugs, or certain prescription medications (e.g.\ thalidomide, isotretinoin, sodium valproate).
    \item \textbf{Maternal infections (TORCH panel):} Intrauterine exposure to infectious pathogens, including Toxoplasmosis, Rubella, Cytomegalovirus, Herpes simplex virus, and Syphilis, which disrupt normal embryogenesis.
    \item \textbf{Nutritional deficiencies:} Inadequate maternal intake of essential micronutrients, most notably folate deficiency, which impairs neural tube closure during early embryogenesis.
    \item \textbf{Maternal chronic conditions:} Poorly controlled maternal diabetes mellitus increases the risk of cardiac and neural anomalies; maternal obesity and thyroid disease are also associated with elevated risk.
    \item \textbf{Physical intrauterine factors:} Mechanical forces that restrict fetal movement or compress the developing fetus, leading to deformations (e.g.\ clubfoot or developmental dysplasia of the hip due to oligohydramnios).
\end{itemize}

\section*{Clinical Features}
The clinical presentation of congenital anomalies depends on the organ systems involved and the severity of the defect.
\begin{itemize}
    \item \textbf{Structural malformations:} Visibly apparent anomalies at birth, such as cleft lip or palate, polydactyly, syndactyly, clubfoot, or microcephaly.
    \item \textbf{Cardiovascular signs:} Cyanosis, heart murmurs, tachypnoea, poor feeding, and failure to thrive in infants with congenital heart disease (e.g.\ Tetralogy of Fallot, ventricular septal defects).
    \item \textbf{Neurological deficits:} Paralysis, sensory loss, neurogenic bladder, or hydrocephalus in babies with neural tube defects like spina bifida.
    \item \textbf{Gastrointestinal obstruction signs:} Bilious vomiting, abdominal distension, and failure to pass meconium within the first 24 to 48 hours of life (suggestive of anorectal malformations, Hirschsprung's disease, or duodenal atresia).
    \item \textbf{Metabolic and endocrine features:} Jaundice, lethargy, poor feeding, or abnormal odour in infants with inborn errors of metabolism (detected via newborn screening).
\end{itemize}

\section*{Differential Diagnosis}
It is important to differentiate congenital anomalies from acquired neonatal conditions and physiological variations:
\begin{itemize}
    \item \textbf{Acquired neonatal infections:} Sepsis, meningitis, or pneumonia, which present with lethargy, poor feeding, and respiratory distress, mimicking cardiac or metabolic congenital disorders.
    \item \textbf{Birth trauma injuries:} Brachial plexus palsy (Erb's palsy), clavicle fracture, or caput succedaneum/cephalohematoma, which present as structural or functional deficits at birth but are acquired during delivery. Distinction: History of difficult delivery and resolution over time.
    \item \textbf{Transient physiological variations:} Physiological jaundice of the newborn, transient murmurs (e.g.\ closing ductus arteriosus), or positional clubfoot that resolves with stretching.
    \item \textbf{Acquired respiratory distress syndrome:} Due to prematurity and surfactant deficiency, rather than structural lung or diaphragmatic anomalies.
    \item \textbf{Maternal drug withdrawal (Neonatal Abstinence Syndrome):} Presents with irritability, tremors, and feeding difficulties, which must be distinguished from metabolic anomalies.
\end{itemize}

\section*{When to Seek Medical Care}
Many major congenital anomalies are detected during routine prenatal scans or newborn examinations. Parents should seek immediate medical attention if a newborn shows signs of physiological distress at home.

\textbf{Call 000 (or local emergency services) for:}
\begin{itemize}
    \item Cyanosis (blue or grey discolouration of the lips, tongue, or face).
    \item Severe respiratory distress: grunting, rapid breathing, or deep chest retractions.
    \item Persistent bilious (green-coloured) vomiting or inability to keep down any fluids.
    \item Extreme lethargy, difficulty waking the baby, or unresponsive behaviour.
    \item Seizures, fits, or abnormal repetitive movements.
\end{itemize}

\section*{Diagnosis and Investigations}
Diagnosis involves a combination of prenatal screening, postnatal clinical examination, and targeted diagnostic testing.
\begin{itemize}
    \item \textbf{Prenatal screening and ultrasound:} First-trimester combined screening (ultrasound nuchal translucency and maternal serum markers), non-invasive prenatal testing (NIPT) via cell-free fetal DNA, and mid-trimester morphology scan (around 18 to 20 weeks).
    \item \textbf{Newborn screening test (Heel prick test):} Performed 48 to 72 hours after birth to screen for cystic fibrosis, congenital hypothyroidism, phenylketonuria (PKU), and other metabolic disorders.
    \item \textbf{Karyotyping and chromosomal microarray:} Performed on amniotic fluid, chorionic villi, or postnatal blood samples to identify chromosomal abnormalities.
    \item \textbf{Echocardiography:} The gold standard for diagnosing structural congenital heart defects in utero or postnatally.
    \item \textbf{Magnetic Resonance Imaging (MRI):} Fetal or neonatal MRI to evaluate complex brain, spine, or thoracic anomalies.
\end{itemize}

\section*{Management}
Management is highly complex and requires a coordinated, multidisciplinary team including paediatricians, surgeons, geneticists, therapists, and support services.
\begin{itemize}
    \item \textbf{Surgical intervention:} Urgent neonatal surgery for life-threatening defects (e.g.\ diaphragmatic hernia, gastroschisis, transposition of the great arteries) or elective surgery for structural reconstruction (e.g.\ cleft lip repair, orthopaedic surgeries).
    \item \textbf{Pharmacological management:} Administration of prostaglandin $E_{1}$ to maintain ductus arteriosus patency in duct-dependent cardiac lesions, or hormone replacement for congenital hypothyroidism.
    \item \textbf{Nutritional support:} Specialised feeding techniques, fortifiers, or enteral tube feeding for infants with feeding difficulties (e.g.\ cleft palate or cardiac failure).
    \item \textbf{Allied health therapies:} Early intervention programmes incorporating physiotherapy, occupational therapy, and speech therapy to optimize motor and cognitive development.
    \item \textbf{Genetic counselling:} Provided to parents to explain the diagnosis, discuss recurrence risks in future pregnancies, and explore family planning options.
    \item \textbf{Psychosocial support:} Counselling and peer support groups for families to assist with the emotional and financial impact of caring for a child with a chronic anomaly.
\end{itemize}

\section*{Complications}
\begin{itemize}
    \item \textbf{Developmental delay and intellectual disability:} Common in children with chromosomal abnormalities or severe central nervous system defects.
    \item \textbf{Chronic organ failure:} Chronic kidney disease (due to renal dysplasia), heart failure (due to uncorrected cardiac defects), or respiratory insufficiency.
    \item \textbf{Growth failure and malnutrition:} Due to high metabolic demands (e.g.\ in congenital heart disease) or feeding difficulties.
    \item \textbf{Recurrent infections:} Increased susceptibility to pneumonia or urinary tract infections (e.g.\ in vesicoureteral reflux).
    \item \textbf{Social isolation and mental health issues:} Elevated risk of anxiety, depression, and social difficulties in both the affected child and their caregivers.
\end{itemize}

\section*{Prognosis and Outcomes}
The prognosis for children with congenital anomalies varies widely depending on the type and severity of the defect, the presence of associated anomalies, and the timing of diagnosis and treatment. Early detection, prenatal planning, and access to tertiary paediatric surgical centres have dramatically improved survival and quality of life over recent decades. While some major anomalies remain associated with high mortality or lifelong disability, many structural differences can be corrected surgically, allowing individuals to lead active, fulfilling lives.

\section*{Prevention and Self-Care}
\begin{itemize}
    \item \textbf{Folate supplementation:} Consuming 400 micrograms of folic acid daily starting at least one month before conception and throughout the first trimester to prevent neural tube defects.
    \item \textbf{Avoid teratogens:} Abstaining from alcohol, tobacco, and recreational drugs during pregnancy, and reviewing all prescription medications with a doctor prior to conceiving.
    \item \textbf{Optimise chronic health conditions:} Achieving good glycaemic control in women with pre-existing diabetes before and during pregnancy.
    \item \textbf{Ensure vaccination is up to date:} Verifying immunity to rubella and varicella prior to pregnancy, and receiving recommended vaccines during pregnancy (e.g.\ pertussis, influenza).
    \item \textbf{Attend prenatal appointments:} Participating in all recommended prenatal check-ups and screening tests to monitor fetal development.
\end{itemize}

\section*{Sources and Further Reading}
The following reputable organisations provide information on congenital anomalies and birth differences:
\begin{itemize}
    \item Australian Institute of Health and Welfare (AIHW). \textit{National monitoring of congenital anomalies in Australia.} https://www.aihw.gov.au/
    \item Healthdirect Australia. \textit{Understanding birth defects and congenital anomalies.} https://www.healthdirect.gov.au/
    \item World Health Organization. \textit{Congenital anomalies factsheets and global prevention strategies.} https://www.who.int/
    \item Centers for Disease Control and Prevention (CDC). \textit{Birth defects research, tracking, and prevention.} https://www.cdc.gov/
    \item Mayo Clinic. \textit{Overview of congenital heart defects and structural anomalies.} https://www.mayoclinic.org/
\end{itemize}
